% Prospective framing only. No final-program outcomes have been read or inserted.
% Existing bibliography keys refer to english/main.tex and RELATED_WORK_AUDIT.md.
\section{Introduction}
Speech emotion recognition (SER) is often evaluated for speakers absent from training. Speaker separation across training, validation and test is an established concern: SERAB and Antoniou et al. explicitly use speaker-exclusive validation as well as testing~\cite{serab,antoniou}. More generally, optimizing a finite-sample selection criterion can affect subsequent performance evaluation~\cite{cawley}. However, a validation score difference and a difference in the selected models' performance on an identical unfamiliar-speaker test are distinct quantities. A larger validation--test gap alone does not establish that the corresponding selection rule produces a worse test model.

Prior SER studies already compare random and actor-independent evaluation across datasets, models or repeated runs~\cite{ibrahim,zielonka,kumari}; matching training counts across speaker/text conditions also has precedent~\cite{atmaja}. Such comparisons establish the relevance of evaluation protocols, but may change training or test membership along with the boundary of interest. We therefore study a specified checkpoint-selection comparison: validation on familiar versus unfamiliar speakers, each using either cross-entropy or unweighted average recall (UAR), along one shared training trajectory and on one shared speaker-exclusive outer test. Holding the candidate trajectory fixed removes differences between training runs from this within-trajectory comparison. The validation groups still contain different people, so the contrast concerns the complete selection policies rather than an isolated voiceprint mechanism.

This design follows earlier analyses of corpora already used in this project~\cite{artifact}. It fixes one model configuration, all 15 candidate epochs and the analysis before the new fits. The principal quantities are the common-test difference between the two cross-entropy policies and its interaction with the choice of cross-entropy versus UAR. All three native corpus tasks and all six planned tests are retained. Full trajectories and a fixed-last-epoch training-group exchange provide descriptive context; they do not supply additional hypothesis tests or determine which results are reported.

% Insert interpretation of the accepted results before this limitations text.
\subsection{Scope of the comparison}
The estimates concern these finite acted-speech corpora and the specified randomization and training program. Repeated draws quantify variation within that program, not sampling of independent languages or corpora. The three tasks differ in their native classes, speaker counts, text coverage and training sizes; their differences cannot identify a language effect. Validation identities and recording distributions also differ between the familiar and unfamiliar groups. Fixing the trajectory therefore does not make speaker exposure an individually randomized causal intervention.

The training-group exchange changes an entire fit population. Its queries also participate in checkpoint selection for the main trajectories, although the exchange descriptor uses the fixed last epoch. It is not a separate globally unused blind test. Likewise, fixed report-group batches preserve padding context without removing the encoder's sensitivity to padding. Neither an interval crossing zero nor the absence of a significant adjusted test establishes equivalence; the 1-pp practical reference is not an equivalence margin tested by this program. Conclusions about which selection policy is preferable must follow the common-test estimates and their uncertainty, including any disagreement across corpora.
